\usepackage{
	amsmath,amssymb,graphicx,booktabs,array,siunitx,
	physics,mathtools,caption,subcaption,float,
	fancyhdr,lastpage,adjustbox,longtable,bm,diagbox,
	makecell,multirow,braket,chngcntr,import,
	tcolorbox,placeins,needspace,etoolbox,environ
}
\tcbuselibrary{breakable,skins,theorems}

\usepackage[top=1.7cm,bottom=1.5cm,left=1.5cm,right=1.8cm]{geometry}
\usepackage[unicode]{hyperref}
\usepackage{bookmark}

\hypersetup{
	colorlinks=true,
	linkcolor=blue,
	citecolor=blue,
	filecolor=magenta,
	urlcolor=blue,
	hidelinks=true
}

% ===== 编号与目录 =====
\counterwithout{section}{chapter}
\renewcommand{\thechapter}{\chinese{chapter}}
\renewcommand{\thesection}{\arabic{section}}

\numberwithin{equation}{section}
\numberwithin{figure}{section}
\numberwithin{table}{section}

\renewcommand{\thefigure}{\thesection.\arabic{figure}}
\renewcommand{\thetable}{\thesection.\arabic{table}}

\setcounter{tocdepth}{2}
\setcounter{secnumdepth}{3}

% ===== 页面与排版 =====
\raggedbottom
\allowdisplaybreaks[3]
\emergencystretch=2em
\setlength{\headheight}{13pt}
\setlength{\parindent}{2em}

% ===== 章节标题页 =====
\makeatletter
\def\@makechapterhead#1{%
	\thispagestyle{fancy}%
	\vspace*{\fill}%
	{\centering\Huge\bfseries 第\thechapter 章\par\vspace{1.0em}#1\par}%
	\vspace*{\fill}%
	\clearpage
}
\def\@makeschapterhead#1{%
	\thispagestyle{fancy}%
	\vspace*{\fill}%
	{\centering\Huge\bfseries #1\par}%
	\vspace*{\fill}%
	\clearpage
}
\makeatother

% ===== 页眉页脚 =====
\pagestyle{fancy}
\fancyhf{}
\fancyhead[LE,RO]{\thepage/\pageref{LastPage}}
\fancyhead[RE]{\nouppercase{\leftmark}}
\fancyhead[LO]{\nouppercase{\rightmark}}
\renewcommand{\headrulewidth}{0.4pt}

\renewcommand{\chaptermark}[1]{\markboth{第\thechapter 章\ #1}{}}
\renewcommand{\sectionmark}[1]{\markright{\thesection\ #1}}

\fancypagestyle{plain}{
	\fancyhf{}
	\fancyhead[LE,RO]{\thepage/\pageref{LastPage}}
	\fancyhead[RE]{\nouppercase{\leftmark}}
	\fancyhead[LO]{\nouppercase{\rightmark}}
	\renewcommand{\headrulewidth}{0.4pt}
}

% ===== 浮动体 =====
\floatplacement{figure}{htbp}
\floatplacement{table}{htbp}

\captionsetup[figure]{labelfont=bf}
\captionsetup[table]{labelfont=bf}

\setcounter{topnumber}{5}
\setcounter{bottomnumber}{3}
\setcounter{totalnumber}{8}

\renewcommand{\topfraction}{0.92}
\renewcommand{\bottomfraction}{0.82}
\renewcommand{\textfraction}{0.06}
\renewcommand{\floatpagefraction}{0.78}

\setlength{\textfloatsep}{10pt plus 2pt minus 3pt}
\setlength{\floatsep}{8pt plus 2pt minus 2pt}
\setlength{\intextsep}{8pt plus 2pt minus 2pt}

\interdisplaylinepenalty=100
\predisplaypenalty=0
\postdisplaypenalty=0

\pretocmd{\section}{\FloatBarrier\Needspace{10\baselineskip}}{}{}
\pretocmd{\subsection}{\Needspace{7\baselineskip}}{}{}

% ===== equation* 自动加 split =====
\makeatletter
\expandafter\let\expandafter\orig@equationstar\csname equation*\endcsname
\expandafter\let\expandafter\orig@endequationstar\csname endequation*\endcsname
\RenewEnviron{equation*}{%
	\orig@equationstar
	\begin{split}
		\BODY
	\end{split}
	\orig@endequationstar
}
\renewcommand{\[}{\orig@equationstar}
\renewcommand{\]}{\orig@endequationstar}
\makeatother

% ===== tcolorbox 通用样式 =====
\tcbset{
	mybox/.style={
		enhanced,
		breakable,
		sharp corners,
		boxrule=0.1pt,
		width=\linewidth,
		enlarge left by=14pt,
		enlarge right by=14pt,
		boxsep=0pt,
		left=12pt,right=12pt,
		top=10pt,bottom=8pt,
		before skip=8pt,
		after skip=8pt,
		detach title,
		fonttitle=\bfseries,
		boxed title style={
			sharp corners,
			boxrule=0.35pt,
			left=8pt,right=8pt,
			top=2.5pt,bottom=2.5pt
		},
		attach boxed title to top left={xshift=10pt,yshift*=-\tcboxedtitleheight/2},
		before upper={
			\setlength{\parindent}{2em}%
			\setlength{\parskip}{0pt}%
			\indent\ignorespaces
		}
	}
}

\newtcolorbox[auto counter,number within=section]{example}[1][]{%
	mybox,
	colback=green!3,
	colframe=green!45!black,
	title={例\ \thetcbcounter},
	boxed title style={colback=green!60!black,colframe=green!60!black},
	#1
}

\newtcolorbox[auto counter,number within=section]{notes}[1][]{%
	mybox,
	colback=blue!3,
	colframe=blue!45!black,
	title={注\ \thetcbcounter},
	boxed title style={colback=blue!60!black,colframe=blue!60!black},
	#1
}

% ===== 数学命令 =====
\ExplSyntaxOn
\cs_if_exist:NTF \vb
	{\RenewDocumentCommand{\vb}{m}{\bm{#1}}}
	{\NewDocumentCommand{\vb}{m}{\bm{#1}}}

\cs_if_exist:NTF \uv
	{\RenewDocumentCommand{\uv}{m}{\hat{\bm{#1}}}}
	{\NewDocumentCommand{\uv}{m}{\hat{\bm{#1}}}}

\cs_if_exist:NTF \d
	{\RenewDocumentCommand{\d}{}{\mathop{}\!\mathrm{d}}}
	{\NewDocumentCommand{\d}{}{\mathop{}\!\mathrm{d}}}
\ExplSyntaxOff

\DeclareMathOperator{\diag}{diag}
\DeclareMathOperator{\sinc}{sinc}
\DeclareMathOperator{\arsinh}{arsinh}
\DeclareMathOperator{\arcsinh}{arsinh}
\DeclareMathOperator{\arcosh}{arcosh}
\DeclareMathOperator{\artanh}{artanh}
\DeclareMathOperator{\sn}{sn}
\DeclareMathOperator{\cn}{cn}
\DeclareMathOperator{\dn}{dn}

\newcommand{\crossdot}{%
	\mathbin{\vcenter{\offinterlineskip
		\hbox to 6pt{\hss$\times$\hss}
		\hbox to 6pt{\hss$\cdot$\hss}}}
}
\newcommand{\dotdotdot}{%
	\mathbin{\vcenter{\offinterlineskip
		\hbox to 4pt{\hss$\cdot$\hss}
		\hbox to 4pt{\hss$\cdot$\hss}
		\hbox to 4pt{\hss$\cdot$\hss}}}
}
